\documentclass[a4paper,11pt]{article}
\usepackage[utf8]{inputenc}
\usepackage[T1]{fontenc}
\usepackage{tgheros}
\renewcommand{\familydefault}{\sfdefault}
\usepackage[english]{babel}
\usepackage{geometry}
\geometry{left=1in, right=1in, top=1in, bottom=1in}
\usepackage{setspace}
\onehalfspacing
\usepackage{parskip}
\usepackage{enumitem}
\usepackage{booktabs}
\usepackage{graphicx}
\usepackage{amsmath}
\usepackage{natbib}
\bibliographystyle{apalike}
\usepackage[colorlinks=true, linkcolor=black, urlcolor=blue, citecolor=black]{hyperref}

\usepackage{titlesec}
\titleformat{\section}{\large\bfseries}{\thesection.}{0.5em}{}
\titleformat{\subsection}{\normalsize\bfseries}{\thesubsection}{0.5em}{}

\begin{document}

%================ TITLE PAGE ================%
\begin{center}
\vspace*{1cm}
{\LARGE \textbf{Equal Opportunity by Design?}} \\[8pt]
{\Large \textbf{The Causal Effects of Higher Education Access Policies}} \\[4pt]
{\Large \textbf{on Labor Market Inequality in Brazil}} \\[24pt]
{\large PhD Research Proposal} \\[16pt]
{\large Baha Khammari} \\[4pt]
M.Sc.\ Economics, University of Cologne \\
\href{mailto:b.e.khammari@gmail.com}{b.e.khammari@gmail.com} \\[4pt]
\today
\end{center}

\vspace{24pt}

%================ EXECUTIVE SUMMARY ================%
\section{Introduction}

Over the past two decades, Brazil has implemented three large-scale policies to expand access to higher education: \textit{ProUni} (2005), a merit-based scholarship programme for low-income students at private institutions; \textit{FIES} (reformed 2010), a subsidised student loan scheme; and the \textit{Quota Law} (Lei de Cotas, 2012), which reserves places at federal universities for public-school graduates, low-income students, and racial minorities. Together, these policies reshaped the composition of Brazil's tertiary-educated workforce. Yet their comparative labour-market effects---and whether they reduce or reproduce existing inequalities---remain poorly understood.

This proposal outlines a three-paper PhD project that provides the first systematic causal comparison of these policies' effects on formal-sector labour market outcomes. Building on my master's thesis, which estimated the wage effects of ProUni using heterogeneity-robust staggered difference-in-differences \citep{callaway2021}, the project extends the analysis in three directions: (i)~comparing ProUni against the Quota Law and FIES to identify which policy design is most effective at reducing inequality, (ii)~examining outcomes beyond wages, including formality transitions, occupational upgrading, and job stability, and (iii)~developing a dose--response framework to inform optimal policy calibration.

The project exploits Brazil's uniquely rich administrative data infrastructure---the RAIS employer--employee registry, ProUni and FIES programme records, and Census microdata---accessible through the \texttt{basedosdados} BigQuery interface, with which I have extensive experience. Methodologically, the project contributes to the rapidly evolving literature on heterogeneous treatment effects in staggered adoption settings \citep{roth2023, dechaisemartin2020, goodmanbacon2021}.

\subsection{Motivation and relevance}

The question of whether higher education expansion translates into better labour market outcomes---and for whom---is among the most policy-relevant in development economics. In Brazil, tertiary enrolment more than doubled between 2003 and 2018, driven primarily by the three policies studied here. This expansion occurred alongside a historic reduction in income inequality (the Gini coefficient fell from 0.59 to 0.52 between 2001 and 2014), but the contribution of education policy to this trend is causally unidentified.

My thesis findings motivate several of the core questions: ProUni's wage effects are non-monotonic (+3.04\% in low-dose regions, near-zero in high-dose regions), gains accrue disproportionately to non-white workers, and na\"ive two-way fixed effects (TWFE) estimation masks these patterns entirely. These findings raise three questions that this PhD project addresses: Do the Quota Law and FIES produce similarly heterogeneous effects? Are scholarship-based and quota-based mechanisms substitutes or complements? And can we identify an optimal ``dosage'' that maximises labour market returns for disadvantaged groups?

%================ BACKGROUND ================%
\section{Background and Scientific Basis}

\subsection{Institutional context}

\textbf{ProUni} (\textit{Programa Universidade para Todos}, Lei n\textsuperscript{o} 11.096/2005) provides full or partial tuition scholarships at private higher education institutions to students from low-income households (family income $\leq 1.5$ or $\leq 3$ minimum wages per capita, respectively) who attended public secondary schools or received full private-school scholarships. Institutions participate voluntarily in exchange for tax exemptions. By 2019, ProUni had awarded over 2.4 million scholarships.

\textbf{FIES} (\textit{Fundo de Financiamento Estudantil}) is a federal student loan programme for private-institution students. Following a major expansion in 2010, FIES disbursed over 2.1 million loans by 2014, with subsidised interest rates and income-contingent repayment. A fiscal consolidation in 2015 sharply reduced new contracts, creating temporal variation useful for causal identification.

\textbf{The Quota Law} (Lei n\textsuperscript{o} 12.711/2012) mandates that federal universities reserve at least 50\% of places for graduates of public secondary schools, with sub-quotas for low-income students and ethnic minorities (Black, Brown/\textit{pardo}, and Indigenous). Implementation was staggered: universities had until 2016 to reach full compliance, creating the variation in treatment timing required for difference-in-differences estimation.

\subsection{Gaps in the literature}

The existing literature evaluates each policy in isolation. Studies of ProUni have focused on enrolment effects \citep{cataldo2020} and institutional quality \citep{aprile2009}, with limited attention to downstream labour market outcomes. Research on the Quota Law has examined academic performance and graduation rates at the institution level but rarely traces quota beneficiaries into the formal labour market. FIES has received attention primarily through its fiscal costs and default rates.

No study to date compares the labour-market returns across these three policies within a unified empirical framework. This gap is consequential: the policies target overlapping but distinct populations, operate through different mechanisms (price subsidies, credit, quantity mandates), and interact in ways that single-policy evaluations cannot capture.

Methodologically, most existing evaluations rely on conventional TWFE estimators, which are biased under staggered adoption with heterogeneous treatment effects \citep{goodmanbacon2021, dechaisemartin2020}. My thesis demonstrated this empirically for ProUni: the TWFE estimate of $-0.65\%$ on wages reverses to $+3.04\%$ in low-dose regions once heterogeneity-robust methods are applied. Similar biases likely affect existing estimates for FIES and the Quota Law.

%================ RESEARCH QUESTIONS ================%
\section{Research Questions and Challenges}

\subsection{Research questions}

The project is organised around three self-contained papers, each addressing a distinct research question:

\textbf{Paper~1: Scholarships vs.\ Quotas---Comparative Labour Market Returns.}
\begin{quote}
\textit{Do ProUni scholarships and Quota Law reservations produce different effects on formal-sector wages, employment, and occupational attainment, and do these effects vary by gender, ethnicity, and region?}
\end{quote}
This paper extends my thesis framework to the Quota Law and provides the first head-to-head comparison. It exploits the staggered rollout of quota compliance across federal universities (2013--2016) and the regional variation in ProUni scholarship allocation.

\textbf{Paper~2: Beyond Wages---Formality, Occupational Upgrading, and Job Stability.}
\begin{quote}
\textit{Do higher education access policies affect labour market outcomes beyond wages, including formal-sector entry, occupational mobility, and employment duration?}
\end{quote}
This paper broadens the outcome set. Brazil's large informal sector (approximately 40\% of employment) means that wage effects in the formal sector capture only part of the return to education. RAIS records allow tracking of formal-sector entry, job tenure, occupational codes (CBO), and employer characteristics over time.

\textbf{Paper~3: Dose--Response and Optimal Policy Design.}
\begin{quote}
\textit{What is the shape of the dose--response relationship between higher education access intensity and labour market outcomes, and what does it imply for optimal policy calibration?}
\end{quote}
This paper builds on my thesis finding of non-monotonic ProUni effects. It formalises the dose--response estimation across all three policies using continuous treatment frameworks \citep{callaway2024} and asks whether there is an identifiable ``saturation point'' beyond which additional access yields diminishing or zero returns.

\subsection{Anticipated challenges}

\begin{itemize}[nosep]
    \item \textbf{Selection into treatment.} ProUni and FIES are voluntary; the Quota Law mandates institutional behaviour but not individual take-up. Each paper must address selection at the appropriate level (microregion, institution, or individual) using doubly-robust methods and sensitivity analysis.
    \item \textbf{Concurrent policies.} The three programmes overlap temporally and geographically. Isolating each policy's contribution requires careful control for the others, leveraging differences in eligibility criteria and rollout timing.
    \item \textbf{Data linkage.} Individual-level linkage across programme records and RAIS is not feasible without CPF identifiers. The project therefore operates at the microregion--year level, following the design of my thesis, supplemented by institution-level analysis where data permit.
    \item \textbf{RAIS coverage.} RAIS captures only formal employment. Effects on informality must be inferred indirectly, e.g.\ through PNAD Cont\'inua survey data or changes in the formal employment share at the microregion level.
\end{itemize}

%================ METHODOLOGY ================%
\section{Methodology}

\subsection{Data sources}

All data are publicly accessible through the \texttt{basedosdados} BigQuery interface, which I used extensively in my thesis to construct a microregion--year panel of 558 regions over 2005--2019.

\begin{table}[h]
\centering
\small
\begin{tabular}{@{}lll@{}}
\toprule
\textbf{Dataset} & \textbf{Source} & \textbf{Key variables} \\
\midrule
RAIS & Ministry of Labour & Wages, occupation, tenure, sector, employer \\
ProUni & Ministry of Education & Scholarships by institution, type, year \\
FIES & FNDE & Loans by institution, amount, year \\
Quota Law compliance & Federal universities & Reservation shares, implementation year \\
Census / PNAD Cont\'inua & IBGE & Demographics, education, informality \\
\bottomrule
\end{tabular}
\end{table}

\subsection{Identification strategy}

\textbf{Papers 1 and 2} use the \citet{callaway2021} group--time average treatment effect estimator, which avoids the bias from ``forbidden comparisons'' in staggered TWFE by restricting contrasts to not-yet-treated or never-treated units. Treatment is defined at the microregion level as the year in which scholarship/quota exposure first exceeds a threshold, following the approach validated in my thesis. Doubly-robust estimation combines inverse probability weighting with outcome regression to reduce sensitivity to model misspecification.

For the Quota Law specifically, the staggered compliance timeline (2013--2016) across 63 federal universities provides clean variation in treatment timing. I construct a municipality-level measure of quota exposure based on the share of the local student-age population within commuting distance of a treated university.

\textbf{Paper~3} extends the binary treatment framework to continuous treatment intensity using the generalised DiD estimator for continuous treatments \citep{callaway2024}. This allows nonparametric estimation of the dose--response function $\tau(d)$, where $d$ measures cumulative scholarship/quota intensity per capita.

\subsection{Robustness and sensitivity}

Following the protocol established in my thesis, each paper includes:
\begin{itemize}[nosep]
    \item Pre-trend testing and HonestDiD sensitivity analysis \citep{rambachan2023} under violations of parallel trends.
    \item Bacon decomposition \citep{goodmanbacon2021} to diagnose TWFE bias.
    \item Placebo tests using pre-treatment cohorts and untreated regions.
    \item Controls for concurrent policies (e.g., \textit{Bolsa Fam\'ilia}, minimum wage adjustments).
    \item Alternative treatment definitions and dose cutoffs.
\end{itemize}

\subsection{Software}

All estimation is conducted in R, using the \texttt{did} package \citep{callaway2021}, \texttt{fixest} for high-dimensional fixed effects \citep{berge2018}, \texttt{HonestDiD} \citep{rambachan2023}, and \texttt{tidyverse}/\texttt{ggplot2} for data processing and visualisation. Data extraction from BigQuery uses SQL via \texttt{basedosdados}.

%================ EXPECTED CONTRIBUTION ================%
\section{Expected Contribution and Impact}

\subsection{Scientific contribution}

\begin{enumerate}[nosep]
    \item \textbf{First comparative evaluation} of Brazil's three main higher education access policies within a unified causal framework, addressing a significant gap in the development economics and education policy literatures.
    \item \textbf{Methodological demonstration} of heterogeneity-robust DiD methods applied to a real-world staggered policy environment, contributing to the growing body of applied work that implements the theoretical advances of \citet{callaway2021}, \citet{dechaisemartin2020}, and \citet{sun2021}.
    \item \textbf{Expanded outcome analysis} beyond wages to formality, occupational upgrading, and job stability---outcomes that are first-order relevant in developing countries with large informal sectors but rarely examined in the education policy literature.
    \item \textbf{Dose--response estimation} for education policies, providing evidence on optimal programme scale and informing the design of future interventions.
\end{enumerate}

\subsection{Policy impact}

The findings are directly relevant to policymakers in Brazil and other middle-income countries designing higher education access programmes. The Quota Law's ten-year review (mandated for 2022, ongoing) and recurrent debates about FIES's fiscal sustainability create immediate demand for rigorous comparative evidence. The dose--response analysis in Paper~3 can inform decisions about programme scale---a question my thesis flagged (diminishing returns in high-dose regions) but could not fully resolve with a single-policy design.

More broadly, the project contributes to the global debate on whether expanding higher education access through different mechanisms (price subsidies, quantity mandates, credit) produces equitable labour market outcomes or merely shifts rents.

%================ RESEARCH ETHICS ================%
\section{Research Ethics}

This project uses exclusively \textbf{publicly available, de-identified administrative data}. RAIS microdata accessed through \texttt{basedosdados} are aggregated at the establishment or microregion level; individual worker identifiers (CPF) are not available in the public release. ProUni and FIES records are published as aggregate statistics by institution and year. Census and PNAD data are anonymised by IBGE prior to release.

No human subjects are directly involved. The project does not require institutional review board (IRB/ethics committee) approval under standard guidelines for secondary analysis of public administrative data, though I will seek confirmation from the host institution's ethics office upon enrolment.

The research addresses inequality by design: it examines whether affirmative-action and access policies achieve their stated equity objectives. Findings will be presented with care to avoid stigmatising specific demographic groups. Following best practice in the economics of discrimination, I report heterogeneous effects by ethnicity, gender, and region to illuminate policy-relevant disparities, not to attribute outcomes to group characteristics.

All code and replication materials will be made publicly available on GitHub upon publication.

%================ TIMELINE ================%
\section{Timeline}

The proposed timeline assumes a three-year structured PhD programme beginning in October 2026.

\begin{table}[h]
\centering
\small
\begin{tabular}{@{}lp{10.5cm}@{}}
\toprule
\textbf{Period} & \textbf{Milestones} \\
\midrule
\textbf{Year 1} & \\
Months 1--4 & Literature review; coursework (as required by programme). Extend thesis data infrastructure to include FIES and Quota Law records. \\
Months 5--8 & \textbf{Paper 1} estimation: Quota Law treatment assignment, staggered DiD, comparative analysis with ProUni results. \\
Months 9--12 & Paper 1 draft completed. Present at internal seminar. Begin Paper 2 data construction (occupational codes, tenure, formality measures). \\[4pt]
\textbf{Year 2} & \\
Months 13--16 & \textbf{Paper 2} estimation: formality transitions, occupational upgrading, job stability. \\
Months 17--20 & Paper 2 draft completed. Submit Paper 1 to journal. Conference presentations (e.g., EEA, VfS, ESPE, LACEA). \\
Months 21--24 & Begin \textbf{Paper 3}: continuous treatment framework, dose--response estimation. Revise Paper 1 based on feedback. \\[4pt]
\textbf{Year 3} & \\
Months 25--28 & Paper 3 estimation and draft. Submit Paper 2. \\
Months 29--32 & Revisions and integration. Job market preparation. \\
Months 33--36 & Dissertation submission. Job market presentations. \\
\bottomrule
\end{tabular}
\end{table}

%================ REFERENCES ================%
\section{References}

\begingroup
\renewcommand{\section}[2]{}
\begin{thebibliography}{99}

\bibitem[Aprile and Barone, 2009]{aprile2009}
Aprile, M.R. and Barone, R.E.M. (2009).
\newblock Educa\c{c}\~ao superior: pol\'iticas p\'ublicas para inclus\~ao social.
\newblock \textit{Revista @mbienteeduca\c{c}\~ao}, 2(1), 39--55.

\bibitem[Berg\'e, 2018]{berge2018}
Berg\'e, L. (2018).
\newblock Efficient estimation of maximum likelihood models with multiple fixed-effects: the R package fixest.
\newblock \textit{CREA Discussion Paper}, 2018-13.

\bibitem[Callaway and Sant'Anna, 2021]{callaway2021}
Callaway, B. and Sant'Anna, P.H.C. (2021).
\newblock Difference-in-differences with multiple time periods.
\newblock \textit{Journal of Econometrics}, 225(2), 200--230.

\bibitem[Callaway et~al., 2024]{callaway2024}
Callaway, B., Goodman-Bacon, A. and Sant'Anna, P.H.C. (2024).
\newblock Difference-in-differences with a continuous treatment.
\newblock \textit{NBER Working Paper}, No.\ 32117.

\bibitem[Cataldo et~al., 2020]{cataldo2020}
Cataldo, B., Cede\~no, D.K.A. and Di~Maio, M. (2020).
\newblock Scholarships and student outcomes: evidence from Brazil's ProUni.
\newblock \textit{IZA Discussion Paper}, No.\ 13748.

\bibitem[de~Chaisemartin and D'Haultf\oe ille, 2020]{dechaisemartin2020}
de~Chaisemartin, C. and D'Haultf\oe ille, X. (2020).
\newblock Two-way fixed effects estimators with heterogeneous treatment effects.
\newblock \textit{American Economic Review}, 110(9), 2964--2996.

\bibitem[Goodman-Bacon, 2021]{goodmanbacon2021}
Goodman-Bacon, A. (2021).
\newblock Difference-in-differences with variation in treatment timing.
\newblock \textit{Journal of Econometrics}, 225(2), 254--277.

\bibitem[Rambachan and Roth, 2023]{rambachan2023}
Rambachan, A. and Roth, J. (2023).
\newblock A more credible approach to parallel trends.
\newblock \textit{Review of Economic Studies}, 90(5), 2555--2591.

\bibitem[Roth et~al., 2023]{roth2023}
Roth, J., Sant'Anna, P.H.C., Bilinski, A. and Poe, J. (2023).
\newblock What's trending in difference-in-differences? A synthesis of the recent econometrics literature.
\newblock \textit{Journal of Econometrics}, 235(2), 2218--2244.

\bibitem[Sun and Abraham, 2021]{sun2021}
Sun, L. and Abraham, S. (2021).
\newblock Estimating dynamic treatment effects in event studies with heterogeneous treatment effects.
\newblock \textit{Journal of Econometrics}, 225(2), 175--199.

\end{thebibliography}
\endgroup

\end{document}
